\chapter*{\Large \center Abstract}

This project explores how the same response can be interpreted as exhibiting two behaviors with divergent connotations in LLM-generated responses. In particular, the project demonstrates how the same generated response may be perceived as “sycophantic” or “empathic” depending on the psychosocial lens through which it is evaluated. The results suggest that the issue is not isolated to LLM-generated responses and that human-authored ones are susceptible to dual interpretations too.
Experiments across four language models (GPT-oss, GPT-4o, Claude, and Llama) demonstrate significant differences in the traits of LLM-generated responses compared to human-authored ones to the same personal query. The project uses posts and comments from popular subreddits (r/aitah and r/anxiety) to provide empirical evidence. The investigation indicates a strong association between certain traits of LLM responses which may be perceived as instances of “sycophancy” or “empathy”. The project operationalizes sycophantic behaviour with Goffman's Theory of Face and empathic behaviour with Carl Rogers' Person Centered Therapy framework. DeepReflect is the comparative system developed in this project for analyzing human and AI-generated responses to personal queries across multiple paradigms of psychosocial behavior. The framework is designed to be extensible, enabling it to incorporate additional psychosocial paradigms such as the Rokeach Value Survey (RVS) as additional instruments for analysis of the traits exhibited by responses. Additionally, it implements and investigates the effects of a control mechanism by inserting expert human transcripts to elicit improvements in the quality of responses with chain-of-thought (CoT) reasoning.
